\section{Circuit Constructions}
\label{sec:circuits}

The abstract definition~\eqref{eq:stereo-block} becomes concrete
once one specifies how to implement the eigenspace-conditioned
unitaries $U_{\lambda_j}$.  We present three examples of increasing
complexity, each with explicit gate counts.

\subsection{Diagonal Hamiltonians}
\label{subsec:diagonal}

Let $H$ be a diagonal Hamiltonian on $n$ qubits with eigenvalues
$\lambda_j \geq 0$ in the computational basis.  No diagonalization
step is needed: the encoding unitary acts directly on the ancilla,
conditioned on the computational-basis state of the system register.

For each eigenvalue, the encoding state is prepared by a
$Y$-rotation:
\begin{equation}
	R_y(2\gamma_j)\ket{0} = \ket{\lambda_j},
	\qquad \gamma_j = \arctan(1/\lambda_j),
	\label{eq:diagonal-ry}
\end{equation}
since $\cos\gamma_j = \lambda_j/\sqrt{1+\lambda_j^2}$ and
$\sin\gamma_j = 1/\sqrt{1+\lambda_j^2}$.  The stereographic signal
operator in each eigenspace becomes
$S_{\lambda_j} = U_{\lambda_j}\PauliZ = R_y(2\gamma_j)$, a pure
rotation.

The block encoding is a uniformly controlled rotation: for each
computational-basis state $\ket{j}$, apply $R_y(2\gamma_j)$ on the
ancilla.  For an $n$-qubit system ($N = 2^n$ eigenvalues):
\begin{equation}
	\begin{quantikz}[column sep=0.5cm, row sep=0.5cm]
		\lstick{$\ket{0}_a$} & \gate[style={inner sep=4pt}]{R_y(2\gamma_j)} & \qw \\
		\lstick{$\ket{j_{n-1}}$} & \ctrl{-1}\vqw{1} & \qw \\
		\lstick{$\vdots\phantom{0}$} & \vdots & \\
		\lstick{$\ket{j_0}$} & \ctrl{-3} & \qw
	\end{quantikz}
	\label{eq:diagonal-circuit}
\end{equation}

\begin{proposition}[Diagonal Gate Count]
	\label{prop:diagonal-gates}
	The uniformly controlled $R_y$ in~\eqref{eq:diagonal-circuit}
	can be decomposed into $\mathcal{O}(2^n)$ CNOT gates and
	$\mathcal{O}(2^n)$ single-qubit rotations using standard
	decomposition techniques~\cite{NielsenChuang2010}.  The total
	gate count is $\mathcal{O}(N)$ where $N = 2^n$ is the number
	of eigenvalues, matching the complexity of standard diagonal
	block encoding but using a single ancilla qubit.
\end{proposition}

\begin{remark}
	Adding a new eigenvalue to the spectrum does not change the
	encoding of any existing eigenvalue: the rotation angles depend
	only on individual eigenvalues $\lambda_j$, not on the global
	property $\max_j \lambda_j$.  In standard block encoding, by
	contrast, changing $\|H\|$ changes $\alpha$ and thus the
	encoding of \emph{every} eigenvalue.
\end{remark}

\subsection{Pauli-\texorpdfstring{$Z$}{Z} Sum Hamiltonians}
\label{subsec:pauli-z}

When additional algebraic structure is present, the gate count
can be reduced below $\mathcal{O}(N)$.  Consider a Hamiltonian
that is a sum of tensor products of Pauli-$Z$ operators and
identity:
\begin{equation}
	H = \sum_{k=1}^{m} c_k\, Z_{i_k},
	\label{eq:pauli-z-ham}
\end{equation}
where $Z_{i_k}$ denotes $Z$ on qubit $i_k$ (and identity
elsewhere) and $c_k > 0$.  Such Hamiltonians include classical
Ising models and arise as the diagonal part of many physically
relevant systems.

Since each $Z_{i_k}$ has eigenvalues $\pm 1$, the eigenvalues
of~$H$ are
$\lambda_{\mathbf{s}} = c_1 s_1 + \cdots + c_m s_m$ with
$s_k \in \{+1, -1\}$, and the Hamiltonian is already diagonal in
the computational basis.  However, the rotation angle
$\gamma_\mathbf{s} = \arctan(1/\lambda_\mathbf{s})$ depends on
$\lambda_\mathbf{s}$ only through the partial sums
$\sum c_k s_k$, which are determined by $m$ bits rather than $n$
qubits.

\begin{proposition}[Pauli-$Z$ Sum Gate Count]
	\label{prop:pauli-z-gates}
	The stereographic block encoding
	of~\eqref{eq:pauli-z-ham} can be implemented by:
	\begin{enumerate}
		\item Computing the partial sum
		      $\lambda = \sum c_k s_k$ into $m$ ancilla
		      qubits via a reversible adder
		      ($\mathcal{O}(\poly(m))$ gates),
		\item applying a controlled-$R_y$ on the encoding
		      ancilla conditioned on the partial-sum
		      register ($\mathcal{O}(\poly(m))$ gates),
		\item uncomputing the partial sum
		      ($\mathcal{O}(\poly(m))$ gates).
	\end{enumerate}
	The total gate complexity is $\mathcal{O}(\poly(m))$
	rather than $\mathcal{O}(2^n)$, making this efficient
	when $m \ll n$.
\end{proposition}

\begin{remark}[Positivity shift]
	\label{rem:positivity-pauli}
	Equation~\eqref{eq:pauli-z-ham} with $c_k > 0$ has
	eigenvalues in $[-\sum c_k, +\sum c_k]$, which includes
	negative values.  A spectral shift by
	$\lambda_{\mathrm{shift}} = \sum_k c_k$ restores
	positivity.
\end{remark}

\subsection{Heisenberg Model}
\label{subsec:heisenberg}

As a non-trivial example requiring a genuine change of basis,
consider the isotropic Heisenberg model with a longitudinal field
on two qubits:
\begin{equation}
	H = J(X \otimes X + Y \otimes Y + Z \otimes Z) + h(Z \otimes I),
	\label{eq:heisenberg}
\end{equation}
where $J$ is the exchange coupling and $h$ the field strength.

The Hamiltonian is block-diagonal in the standard basis:
$\ket{00}$ and $\ket{11}$ are eigenstates with eigenvalues
$\lambda_1 = J + h$ and $\lambda_4 = J - h$, while the
$\{\ket{01}, \ket{10}\}$ subspace has the $2\times 2$ block
\begin{equation}
	\begin{pmatrix} -J+h & 2J \\ 2J & -J-h \end{pmatrix}
\end{equation}
with eigenvalues
\begin{equation}
	\lambda_{2,3} = -J \pm \sqrt{h^2 + 4J^2}.
	\label{eq:heisenberg-spectrum}
\end{equation}

The change-of-basis unitary $V$ that diagonalizes $H$ acts
non-trivially only on the $\{\ket{01}, \ket{10}\}$ subspace as
a Givens rotation.  Explicitly,
\begin{equation}
	V = \begin{pmatrix}
		1 & 0           & 0            & 0 \\
		0 & \cos\varphi & -\sin\varphi & 0 \\
		0 & \sin\varphi & \cos\varphi  & 0 \\
		0 & 0           & 0            & 1
	\end{pmatrix},
	\label{eq:givens-rotation}
\end{equation}
where $\tan(2\varphi) = 2J/h$ (or $\varphi = \pi/4$ when
$h = 0$, giving Bell-state eigenvectors).  This decomposes into
$\mathcal{O}(1)$ elementary gates (a controlled-$R_y$ and
CNOT gates).

The full stereographic block encoding on 3 qubits (2 system +
1 ancilla) is
\begin{equation}
	U = (I_a \otimes V^\dagger) \cdot U_{\mathrm{diag}}
	\cdot (I_a \otimes V),
	\label{eq:heisenberg-stereo}
\end{equation}
where $U_{\mathrm{diag}}$ is the diagonal stereographic
encoding:
\begin{equation}
	\begin{quantikz}[column sep=0.4cm, row sep=0.5cm]
		\lstick{$\ket{0}_a$} & \qw & \gate[style={inner sep=4pt}]{R_y(2\gamma_j)} & \qw & \qw \\
		\lstick{$\ket{\psi_1}$} & \gate[2]{V} & \ctrl{-1}\vqw{1} & \gate[2]{V^\dagger} & \qw \\
		\lstick{$\ket{\psi_0}$} & \qw & \ctrl{-1} & \qw & \qw
	\end{quantikz}
	\label{eq:heisenberg-circuit}
\end{equation}

\begin{proposition}[Heisenberg Gate Count]
	\label{prop:heisenberg-gates}
	The total gate count for the stereographic block encoding
	of the two-qubit Heisenberg model~\eqref{eq:heisenberg} is
	$\mathcal{O}(1)$: a constant number of elementary gates
	independent of $J$, $h$, and $\|H\|$.  Specifically:
	\begin{itemize}
		\item The Givens rotation $V$ decomposes into 2 CNOT
		      gates and 1 controlled-$R_y$ (the controlled-$R_y$
		      itself costs 2 CNOTs and 2 single-qubit
		      rotations~\cite{NielsenChuang2010}), for a total
		      of 4 CNOTs and 2 single-qubit gates.
		\item The uniformly controlled $R_y(2\gamma_j)$ on
		      the ancilla, conditioned on the 2-qubit system
		      register, decomposes into at most 4 CNOTs and 4
		      single-qubit rotations (using the recursive
		      decomposition for uniformly controlled rotations).
		\item The undiagonalization $V^\dagger$ has the same
		      cost as $V$: 4 CNOTs and 2 single-qubit gates.
	\end{itemize}
	The total is at most $12$ CNOTs and $8$ single-qubit
	rotations for the block encoding $U$.  Each QSP step adds
	1 phase gate (on the ancilla) and 1 signal operator
	application (the block encoding), giving a total circuit
	depth of $\mathcal{O}(d)$ where $d$ is the QSP degree.
\end{proposition}

\begin{remark}[Positivity shift for Heisenberg]
	\label{rem:positivity-heisenberg}
	For $J > 0$, the smallest eigenvalue
	$\lambda_{\min} = -J - \sqrt{h^2+4J^2}$ is negative.
	Shifting by $|\lambda_{\min}|$ makes all eigenvalues
	non-negative, with the shifted minimum at zero (mapping
	to the south pole of the Bloch sphere under stereographic
	encoding).
\end{remark}

\subsection{Scalability: Three-Qubit Heisenberg Chain}
\label{subsec:three-qubit}

As an indication of scalability, consider the open three-qubit
Heisenberg chain
\begin{equation}
	H_3 = J(\vec{\sigma}_1 \cdot \vec{\sigma}_2
	+ \vec{\sigma}_2 \cdot \vec{\sigma}_3)
	+ h\, Z_1.
	\label{eq:heisenberg-3}
\end{equation}
The Hilbert space dimension is $2^3 = 8$, so stereographic
block encoding uses 4 qubits (3 system + 1 ancilla).  The
key question is whether the diagonalization $V_3$ admits an
efficient circuit implementation.

For the two-qubit case, the Givens rotation has
$\mathcal{O}(1)$ gate count because the Hamiltonian is
block-diagonal with at most $2\times 2$ blocks.
The three-qubit Hamiltonian has larger blocks: the total
$S_z$ is still conserved, giving sectors of dimension 1
($S_z = \pm 3/2$), 3 ($S_z = +1/2$), and 3 ($S_z = -1/2$).
Diagonalizing the $3\times 3$ blocks requires two Givens
rotations per block, giving $\mathcal{O}(1)$ two-qubit
gates---still constant, though with a larger constant than
the two-qubit case.

More generally, for an $n$-qubit Heisenberg chain with
conserved $S_z$, the largest symmetry sector has dimension
$\binom{n}{\lfloor n/2 \rfloor} = \Theta(2^n/\sqrt{n})$,
and the diagonalization cost grows accordingly.  For small
$n$ (up to $\sim\!6$ qubits), the diagonalization is
tractable via exact Givens decompositions.  Beyond this,
the construction requires either additional structure
(e.g., Bethe ansatz circuits~\cite{Bethe1931}) or
approximate diagonalization.
